\section{The Degeneracy}
\label{sec:degeneracy}

This section states the paper's central claim precisely enough to be checked.
The argument is short. Its consequences are not.

One object plays two roles throughout, and keeping them straight removes most of
the apparent complexity. The correctness check $c$ is what an agent runs to
accept its own draft \emph{and} the verifiable reward a post-training loop
maximizes---the same predicate wearing two hats. The optimum-degeneracy result
(Proposition~\ref{prop:degeneracy}) indicts both at once: an agent accepts
correct-but-serial code, and a training optimum sits on it. The signal-collapse
result (Proposition~\ref{prop:collapse}) is specific to the training loop. The
efficiency-gated reward of Section~\ref{sec:reward} is, correspondingly, both the
acceptance gate our critic enforces at inference and the reward we recommend for
post-training. Where a claim applies to only one role we say so.

\subsection{Setup and notation}

Let $\mathcal{T}$ be a set of parallel-programming tasks. Each task
$t \in \mathcal{T}$ is given by a natural-language prompt with an explicit
function contract, and is equipped with a trusted serial reference $R_t$: a
sequential implementation that defines the correct output on any input. A
policy $\pi$ samples a program $p \sim \pi(\cdot \mid t)$.

A parallel program admits a family of execution configurations. We fix
attention on the OpenMP thread count $n$ and write $c(p) \in \{0,1\}$ for the
correctness indicator: $c(p)=1$ iff $p$ compiles, runs, and matches $R_t$ on
every validation input at every thread count and repetition tested. This is
the strongest output-based check in common use---strictly stronger than a
single execution, which is the usual practice~\cite{pareval}---and the argument
below applies to it, and therefore a fortiori to weaker checks.

For a program that compiles and validates, let $T_p(n)$ be its wall time on
$n$ threads and $T_{R_t}$ the wall time of the serial reference. We use three
performance quantities:
\begin{align}
  S_p(n) &= T_{R_t}/T_p(n) && \text{(reference speedup)}\\
  \Sigma_p(n) &= T_p(1)/T_p(n) && \text{(self-speedup, strong scaling)}\\
  Q_p &= T_{R_t}/T_p(1) && \text{(single-thread quality)}
\end{align}
which satisfy the identity
\begin{equation}
  S_p(n) = Q_p \cdot \Sigma_p(n).
  \label{eq:decomposition}
\end{equation}
Equation~\eqref{eq:decomposition} separates speedup obtained by writing a
better sequential algorithm from speedup obtained by using more cores. It will
matter twice: once because a reward that ignores it is gameable, and once
because a reward that uses only $\Sigma$ is gameable differently.

The correctness-only verifiable reward is
\begin{equation}
  \Rcorr(p) = c(p).
  \label{eq:rcorr}
\end{equation}

\subsection{The serial projection}

\begin{definition}[Serial projection]
\label{def:projection}
The serial projection $\serialproj$ maps a program $p$ to the program obtained
by deleting every OpenMP directive from its source: every line beginning
\texttt{\#pragma omp}, and every OpenMP runtime call, replaced by its
single-threaded semantics.
\end{definition}

$\serialproj$ is a purely syntactic transform. It requires no understanding of
the program. It is implementable in a few dozen lines, and we release ours.

\begin{proposition}[Optimum degeneracy]
\label{prop:degeneracy}
Let $p$ be a program with $c(p)=1$ whose parallel regions are free of
side effects outside the memory the corresponding sequential loop would touch.
Then $c(\serialproj(p)) = 1$, and therefore
\begin{equation*}
  \Rcorr(\serialproj(p)) = \Rcorr(p) = \max_{p'} \Rcorr(p').
\end{equation*}
Consequently $\serialproj(p) \in \arg\max \Rcorr$, and
$\Sigma_{\serialproj(p)}(n) \le 1$ for all $n$.
\end{proposition}

\begin{proof}[Argument]
OpenMP worksharing constructs are, by specification, semantically transparent
in the absence of data races: executing a \texttt{parallel for} region with one
thread is equivalent to executing the loop body sequentially in iteration
order. Removing the directive therefore yields the sequential execution that
$c$ compares against $R_t$. If $p$ was correct at $n=1$---which
$c(p)=1$ requires---then $\serialproj(p)$ produces that same output at every
$n$, because it has no parallelism whose schedule could vary. Hence
$c(\serialproj(p))=1$. Since $\Rcorr$ takes values in $\{0,1\}$ and
$\Rcorr(p)=1$ is maximal, $\serialproj(p)$ attains the maximum. It contains no
parallel regions, so it cannot exceed unit self-speedup.
\end{proof}

The proposition is elementary, and we do not present it as difficult. We
present it because the field is currently deploying $\Rcorr$ as a training
objective for exactly this task, and because it has a consequence that is not
elementary: \emph{the correctness reward can be maximized without ever
producing a parallel program, so no procedure that maximizes it can be said to
have learned parallelism.} Any parallelism observed under such an objective is
inherited from the pretrained prior. The reward neither reinforces nor
protects it.

Note also what the proposition does \emph{not} depend on. It does not depend on
the checker being weak. Strengthening $c$---more inputs, more thread counts,
more repetitions, a happens-before race detector alongside---raises the bar for
$p$ and leaves $\serialproj(p)$ untouched, because a program with no
concurrency passes every concurrency check trivially. Improving the correctness
verifier moves the reward's argmax \emph{toward} the serial program, not away
from it. This is the sense in which the degeneracy is structural.

\subsection{Signal collapse}

Proposition~\ref{prop:degeneracy} concerns where the optimum lies. The second
consequence concerns whether any gradient points anywhere at all.

Group-relative policy optimization and its variants estimate an advantage for
each sample in a group of size $k$ by centering on the group mean and
normalizing by the group standard
deviation~\cite{grpo,dapo,drgrpo,gspo,threeops}:
\begin{equation}
  \hat{A}_i = \frac{r_i - \mu(\mathbf{r})}{\varsigma(\mathbf{r})}, \qquad
  \mathbf{r} = (r_1,\dots,r_k).
  \label{eq:advantage}
\end{equation}

\begin{proposition}[Signal collapse]
\label{prop:collapse}
Let $p_t$ be the per-task probability that a sample is correct. Under
$\Rcorr$, the probability that a group of $k$ i.i.d.\ samples for task $t$
yields $\varsigma(\mathbf{r}) = 0$---and hence a zero or undefined
advantage for every member---is
\begin{equation}
  \Pr[\text{dead group}] = p_t^{\,k} + (1-p_t)^{\,k}.
  \label{eq:deadgroup}
\end{equation}
This quantity approaches $1$ as $p_t \to 1$. Under any reward of the form
$\mathbb{1}[c]\cdot f$ with $f$ continuously distributed over correct programs,
the same probability is at most $(1-p_t)^k$.
\end{proposition}

\begin{proof}[Argument]
$\Rcorr$ takes two values, so the group reward vector has zero variance exactly
when all $k$ samples are correct or all $k$ are incorrect, giving
\eqref{eq:deadgroup}. Under $\mathbb{1}[c]\cdot f$, an all-correct group has
zero variance only if all $k$ measured $f$ values coincide, which has
probability zero for a continuous $f$; the all-incorrect case is unchanged.
\end{proof}

The practical reading of Proposition~\ref{prop:collapse} is uncomfortable. The
regime in which a correctness reward stops teaching anything is precisely the
regime of a \emph{capable} model on a \emph{well-specified} benchmark. Effort
spent making models better at correctness, or benchmarks cleaner, drives $p_t$
up and drives the usable signal down. DAPO's dynamic sampling~\cite{dapo}
discards zero-variance groups explicitly, which is the correct engineering
response and also an admission of the problem: on a saturated task
distribution, most of the sampling budget is discarded.

\subsection{Correctness is portable; performance is not}

\begin{observation}[Machine invariance]
\label{obs:invariance}
$\Rcorr$ is invariant to thread count, loop schedule, thread affinity,
compiler, and target architecture. $S_p(n)$ and $\Sigma_p(n)$ are not.
\end{observation}

For most software this is a virtue: a correctness oracle that changed with the
machine would be a bad oracle. For HPC it identifies the mismatch. The purpose
of writing OpenMP rather than a sequential loop is to exploit a particular
machine. A reward that is constant across every machine therefore defines an
optimum that is, by construction, independent of the only thing the task is
about---and an optimum that is tied to whatever machine the timing was collected
on, which is why we measure on a fixed machine and treat cross-machine transfer
as an explicit limitation (Section~\ref{sec:discussion}).

\subsection{Why more verification does not fix it}

It is tempting to read Proposition~\ref{prop:degeneracy} as a call for a better
verifier, and we spent some time making that mistake. The relevant distinction
is between \emph{extensional} verification---does the program produce the right
outputs---and \emph{intensional} verification---does it produce them the right
way~\cite{gamingverifiers}. Differential testing across thread counts,
repetition to expose nondeterminism, and happens-before race
detection~\cite{archer,tsan} are all improvements to the extensional check.
They tighten the set of accepted programs. They do not remove the serial
program from it, because the serial program is extensionally perfect.

What is needed instead is a second observable that is not a function of the
output. Wall time is the obvious one, it is what the HPC community already
measures, and it is the one we adopt in Section~\ref{sec:reward}. Its cost is
that it is a physical measurement rather than a logical predicate, with all
the reproducibility difficulty that implies---which is why
Section~\ref{sec:results-robust} treats measurement reliability as a
first-class result rather than an appendix.
